\documentclass[tikz,border=6pt]{standalone}
\usepackage{ctex}
\usepackage{amsmath}
\usetikzlibrary{calc, arrows.meta, decorations.markings}

\begin{document}
\begin{tikzpicture}[
  scale=1.1, thick,
]

% 坐标轴（参考）
\draw[gray!40, thin, ->] (-0.3, 0) -- (5.5, 0) node[right, font=\footnotesize, gray] {$x$};
\draw[gray!40, thin, ->] (0, -0.3) -- (0, 3.8) node[above, font=\footnotesize, gray] {$y$};

% l1：直线 y = 0.6x + 2.8，x: 0→4.5
\draw[black, thick] (0.0, 2.8) -- (4.5, {2.8 + 0.6*4.5});
\node[font=\small] at (4.7, {2.8 + 0.6*4.5}) {$l_1$};
\node[font=\footnotesize] at (1.6, 3.55) {$Ax+By+C_1=0$};

% l2：直线 y = 0.6x + 0.9，x: 0→4.5
\draw[blue, dashed, thick] (0.0, 0.9) -- (4.5, {0.9 + 0.6*4.5});
\node[font=\small, blue] at (4.7, {0.9 + 0.6*4.5}) {$l_2$};
\node[font=\footnotesize, blue] at (1.6, 0.55) {$Ax+By+C_2=0$};

% 公共垂线段：垂直于两平行线
% 两线斜率 0.6，法线斜率 -1/0.6 = -5/3
% 取 x=2.0 处：l1 上 y = 2.8+1.2=4.0，l2 上 y=0.9+1.2=2.1
% 法线过 l2 点 (2.0, 2.1)，方向 (1, -5/3) 标准化后倒数到 l1
% 法线方程：y - 2.1 = -5/3*(x-2.0)
% 与 l1 交点：y=0.6x+2.8 → 0.6x+2.8 - 2.1 = -5/3*(x-2.0)
% 0.6x + 0.7 = -5/3 x + 10/3
% 0.6x + 5/3 x = 10/3 - 0.7
% x(0.6 + 1.667) = 3.333 - 0.7 = 2.633
% x * 2.267 = 2.633 → x = 1.161
% y = 0.6*1.161 + 2.8 = 0.697 + 2.8 = 3.497
\coordinate (Q1) at (1.161, 3.497);  % l1 上垂足
\coordinate (Q2) at (2.0, 2.1);       % l2 上起点

\draw[red, thick] (Q1) -- (Q2);

% 直角标记（在 Q2 处）
% 法线方向 Q2→Q1: (1.161-2.0, 3.497-2.1) = (-0.839, 1.397)，长 1.626
% l2 方向 (1, 0.6)，单位化除以sqrt(1.36)=1.166
\draw[gray, thin]
  ($(Q2) + 0.18*(0.857, 0.515)$) --
  ($(Q2) + 0.18*(0.857, 0.515) + 0.18*(-0.516, 0.857)$) --
  ($(Q2) + 0.18*(-0.516, 0.857)$);

% 直角标记（在 Q1 处）
\draw[gray, thin]
  ($(Q1) + 0.18*(0.857, 0.515)$) --
  ($(Q1) + 0.18*(0.857, 0.515) + 0.18*(-0.516, 0.857)$) --
  ($(Q1) + 0.18*(-0.516, 0.857)$);

% 距离标注 d
\node[font=\small, red] at (1.25, 2.75) {$d$};

% 端点圆点
\filldraw[black] (Q1) circle (1.5pt);
\filldraw[black] (Q2) circle (1.5pt);

% 公式
\node[font=\footnotesize, align=center] at (3.8, 1.5)
  {$d = \dfrac{|C_1 - C_2|}{\sqrt{A^2+B^2}}$\\[4pt]（需先化同 $A,B$）};

\end{tikzpicture}
\end{document}
